% Grado 3 -- generado por tools/build_tex.py a partir de raw/grado_03.txt
\section{Grado 3}\label{grado:3}

% PDF p. 22
\dba{1}{Interpreta, formula y resuelve problemas aditivos de composición, transformación y comparación en diferentes contextos; y multiplicativos, directos e inversos, en diferentes contextos.}

\evidencias

\begin{itemize}
  \item Construye diagramas para representar las relaciones observadas entre las cantidades presentes en una situación.
  \item Resuelve problemas aditivos (suma o resta) y multiplicativos (multiplicación o división) de composición de medida y de conteo.
  \item Propone estrategias para calcular el número de combinaciones posibles de un conjunto de atributos.
  \item Analiza los resultados ofrecidos por el cálculo matemático e identifica las condiciones bajo las cuales ese resultado es o no plausible.
\end{itemize}

\ejemplo

En los partidos de baloncesto, una cesta puede tener un valor de tres puntos, de dos puntos o de un punto. Propone el número de cestas que hizo cada equipo.

\fig{g03_p22_1}{Marcador de baloncesto: Local 45 – Visitante 62.}

En la imagen se muestra el marcador al finalizar el primer tiempo de un partido de baloncesto.

Si el partido terminó empatado en 80 puntos responde: ¿Cuántas cestas hizo el equipo LOCAL?, ¿Cuántas cestas hizo el equipo VISITANTE? ¿Cómo pudo haberse logrado los puntajes?

Si hubo tres tiros libres para LOCAL y cinco tiros libres para VISITANTE, determina el puntaje de cada equipo sabiendo que cada tiro libre vale un solo punto.

% PDF p. 22
\dba{2}{Propone, desarrolla y justifica estrategias para hacer estimaciones y cálculos con operaciones básicas en la solución de problemas.}

\evidencias

\begin{itemize}
  \item Utiliza las propiedades de las operaciones y del Sistema de Numeración Decimal para justificar acciones como: descomposición de números, completar hasta la decena más cercana, duplicar, cambiar la posición, multiplicar abreviadamente por múltiplos de 10, entre otros.
  \item Reconoce el uso de las operaciones para calcular la medida (compuesta) de diferentes objetos de su entorno.
  \item Argumenta cuáles atributos de los objetos pueden ser medidos mediante la comparación directa con una unidad y cuáles pueden ser calculados con algunas operaciones entre números.
\end{itemize}

\ejemplo

El siguiente gráfico presenta la cantidad de dinero de un juego de mesa que tiene Juan para repartir por partes iguales entre sus tres hijos.

\fig{g03_p23_2}{Tres fichas de 10 y cinco monedas de 1.}

\fig{g03_p23_1}{Cuatro fichas (billetes) de 100.}

Escribe la cantidad de dinero que Juan va a repartir.

Thomas, el hijo mayor de Juan, hizo el siguiente cálculo:

\fig{g03_p23_3}{Hoja con procedimientos escritos de cálculo (algoritmos de suma y resta) de dos estudiantes.}

Describe el procedimiento realizado por Thomas y explica por qué es o no válido. Propone cambios al procedimiento de tal manera que sea válido, en caso de considerarlo incorrecto.

% PDF p. 23
\dba{3}{Establece comparaciones entre cantidades y expresiones que involucran operaciones y relaciones aditivas y multiplicativas y sus representaciones numéricas.}

\evidencias

\begin{itemize}
  \item Realiza mediciones de un mismo objeto con otros de diferente tamaño y establece equivalencias entre ellas.
  \item Utiliza las razones y fracciones como una manera de establecer comparaciones entre dos cantidades.
  \item Propone ejemplos de cantidades que se relacionan entre sí según correspondan a una fracción dada.
  \item Utiliza fracciones para expresar la relación de “el todo” con algunas de sus “partes”, asimismo diferencia este tipo de relación de otras como las relaciones de equivalencia (igualdad) y de orden (mayor que y menor que).
\end{itemize}

\ejemplo

Algunos sistemas económicos usan monedas para representar fracciones de la unidad. En Estados Unidos se utiliza como unidad un dólar, y algunas monedas representan fracciones de esta unidad. La siguiente imagen presenta dos ejemplos de las monedas en este sistema, una corresponde a medio dólar y la otra a un cuarto de dólar.

\fig{g03_p23_4}{Monedas por ambas caras.}

En Colombia, las monedas que representaban fracciones de peso, desaparecieron hace ya varios años. Existían monedas con denominaciones de uno, dos, cinco, diez, veinte, veinticinco y cincuenta centavos. Construye ese tipo de monedas. En una de las caras diseña la denominación de la moneda en centavos, y en la otra, diseña la denominación de la moneda en fracción (semejante a como se denominan en Estados Unidos).

\begin{itemize}
  \item La moneda de un cuarto de peso, ¿a cuántos centavos equivaldría?
  \item La moneda de 10 centavos ¿a cuál fracción equivaldría? Si 50 centavos son lo mismo que dos monedas de 20 centavos y una de 10 centavos. Es decir, 50 = 2 (20) + 1 (10)
\end{itemize}
En fracciones sería: medio de peso equivale a 2 monedas de quinto y una moneda de décimo. Es decir, 1 medio = 2 quintos + 1 décimo

\begin{itemize}
  \item Mónica tiene 70 centavos en monedas, Carlos tiene dos monedas de 20 centavos, Paula tiene cinco monedas de 10 centavos. Representa estos valores usando la denominación en forma de fracción de cada moneda.
  \item Propone otras equivalencias para cantidades diferentes de monedas usando tanto la denominación en forma de fracción como en centavos.
\end{itemize}

% PDF p. 24
\dba{4}{Describe y argumenta posibles relaciones entre los valores del área y el perímetro de figuras planas (especialmente cuadriláteros).}

\evidencias

\begin{itemize}
  \item Toma decisiones sobre la magnitud a medir (área o longitud) según la necesidad de una situación.
  \item Realiza recubrimientos de superficies con diferentes figuras planas.
  \item Mide y calcula el área y el perímetro de un rectángulo y expresa el resultado en unidades apropiadas según el caso.
  \item Explica cómo figuras de igual perímetro pueden tener diferente área.
\end{itemize}

\ejemplo

Determina el número de rectángulos que se pueden formar utilizando 12 palillos y en cada caso encuentra el número de cuadrados, cuyo lado es un palillo; que caben en cada rectángulo.

\fig{g03_p24_1}{Dos regiones planas (siluetas) para comparar su superficie y su contorno.}

Si cada uno de los rectángulos formados se imagina como tablas a las que se pone cinta alrededor, indica la cantidad de cinta que se necesita. Da la medida en términos del número de palillos. Dice en cuál de esos rectángulos se usa mas cinta y cuáles menos.

% PDF p. 25
\dba{5}{Realiza estimaciones y mediciones de volumen, capacidad, longitud, área, peso de objetos o la duración de eventos como parte del proceso para resolver diferentes problemas.}

\evidencias

\begin{itemize}
  \item Compara objetos según su longitud, área, capacidad, volumen, etc.
  \item Hace estimaciones de longitud, área, volumen, peso y tiempo según su necesidad en la situación.
  \item Hace estimaciones de volumen, área y longitud en presencia de los objetos y los instrumentos de medida y en ausencia de ellos.
  \item Empaca objetos en cajas y recipientes variados y calcula la cantidad que podría caber; para ello tiene en cuenta la forma y volumen de los objetos a empacar y la capacidad del recipiente en el que se empaca.
\end{itemize}

\ejemplo

Se tienen que empacar frascos de 8 cm de diámetro y 15 cm de alto. El empacador dispone de cajas de base rectangular de diferentes tamaños y tiene que decidir la caja de tamaño más adecuado. Explica diversos procedimientos que el empacador puede seguir para tomar la decisión más adecuada. Identifica las medidas de tres posibles cajas, si por peso se sugiere que en cada una vayan 50 frascos.

% PDF p. 25
\dba{6}{Describe y representa formas bidimensionales y tridimensionales de acuerdo con las propiedades geométricas.}

\evidencias

\begin{itemize}
  \item Relaciona objetos de su entorno con formas bidimensionales y tridimensionales, nombra y describe sus elementos.
  \item Clasifica y representa formas bidimensionales y tridimensionales tomando en cuenta sus características geométricas comunes y describe el criterio utilizado.
  \item Interpreta, compara y justifica propiedades de formas bidimensionales y tridimensionales.
\end{itemize}

\ejemplo

La profesora de tercero tiene sobre su mesa los cuerpos geométricos que se ven en la imagen:

\fig{g03_p25_1}{Cuerpos geométricos: prisma rectangular, prisma triangular y cilindro.}

\begin{itemize}
  \item David y María no pudieron ver los cuerpos geométricos de la profesora pues no asistieron a clase. Ellos deben realizar la construcción de los mismos con cartulina, cinta y tijeras de tal manera que tengan la misma forma que los de la profesora.
  \item Envía por escrito un mensaje preciso a David y María para que puedan realizar la construcción requerida. El mensaje no puede incluir dibujos, solo las indicaciones adecuadas de tal manera que puedan construir los cuerpos basándose en las indicaciones.
\end{itemize}
Patricia y Román quisieron ayudar a David y María. Para ello escribieron los siguientes mensajes:

Revisa los mensajes escritos e indica si con ellos David y María pueden construir de forma igual los cuerpos que tenía la profesora sobre la mesa y mejora los mensajes escritos.

% PDF p. 26
\dba{7}{Formula y resuelve problemas que se relacionan con la posición, la dirección y el movimiento de objetos en el entorno.}

\evidencias

\begin{itemize}
  \item Localiza objetos o personas a partir de la descripción o representación de una trayectoria y construye representaciones pictóricas para describir sus relaciones.
  \item Identifica y describe patrones de movimiento de figuras bidimensionales que se asocian con transformaciones como: reflexiones, traslaciones y rotaciones de figuras.
  \item Identifica las propiedades de los objetos que se conservan y las que varían cuando se realizan este tipo de transformaciones.
  \item Plantea y resuelve situaciones en las que se requiere analizar las transformaciones de diferentes figuras en el plano.
\end{itemize}

\ejemplo

En un concurso de fotografías Tomás y Alejandro presentan un mosaico con mariposas. Escribe algunas condiciones para que se incluyan una tercera y una cuarta columna de fotografías que conserven la forma como se disponen las imágenes de las dos primeras columnas.

\fig{g03_p26_1}{Secuencia de imágenes de una figura (hoja) girada y reflejada: transformaciones en el plano.}

% PDF p. 27
\dba{8}{Describe y representa los aspectos que cambian y permanecen constantes en secuencias y en otras situaciones de variación.}

\evidencias

\begin{itemize}
  \item Describe de manera cualitativa situaciones de cambio y variación utilizando lenguaje natural, gestos, dibujos y gráficas.
  \item Construye secuencias numéricas y geométricas utilizando propiedades de los números y de las figuras geométricas.
  \item Encuentra y representa generalidades y valida sus hallazgos de acuerdo al contexto.
\end{itemize}

\ejemplo

El gráfico muestra arreglos triangulares de puntos. En la primera posición se tiene 1punto, en la segunda 3 puntos, en la tercera 6 puntos, en la cuarta 10 puntos. Registra (en su orden) el número de puntos en cada posición:

\fig{g03_p27_1}{Arreglos triangulares de puntos: 1, 3, 6, 10, 15, …, 28 (números triangulares).}

\fig{g03_p27_2}{Tabla para completar: posición (primero, segundo, …, noveno), número de puntos (1, 3, 6, …) y descripción del proceso para obtener el siguiente arreglo.}

Explica cómo encontrar el número de puntos en una posición cualquiera. Justifica si existe un arreglo triangular que tenga 35 puntos o 38 puntos.

% PDF p. 27
\dba{9}{Argumenta sobre situaciones numéricas, geométricas y enunciados verbales en los que aparecen datos desconocidos para definir sus posibles valores según el contexto.}

\evidencias

\begin{itemize}
  \item Propone soluciones con base en los datos a pesar de no conocer el número.
  \item Toma decisiones sobre cantidades aunque no conozca exactamente los valores.
  \item Trabaja sobre números desconocidos y con esos números para dar respuestas a los problemas.
\end{itemize}

\ejemplo

José y Patricia tienen cada uno, una caja de dulces. No se sabe cuántos dulces hay en cada caja, pero sí que cada caja tiene la misma cantidad. José tiene un dulce extra encima de su caja. Patricia tiene tres dulces encima de su caja. Dibuja o escribe con letras o con otros símbolos, cuántos dulces tienen entre José y Patricia. ¿Quién tiene más dulces? ¿Cuántos dulces se necesita dar a uno de ellos para que tenga la misma cantidad que el otro? Explica la respuesta y los procedimientos.

% PDF p. 28
\dba{10}{Lee e interpreta información contenida en tablas de frecuencia, gráficos de barras y/o pictogramas con escala, para formular y resolver preguntas de situaciones de su entorno.}

\evidencias

\begin{itemize}
  \item Identifica las características de la población y halla su tamaño a partir de diferentes representaciones estadísticas.
  \item Construye tablas y gráficos que representan los datos a partir de la información dada.
  \item Analiza e interpreta información que ofrecen las tablas y los gráficos de acuerdo con el contexto.
  \item Identifica la moda a partir de datos que se presentan en gráficos y tablas.
  \item Compara la información representada en diferentes tablas y gráficos para formular y responder preguntas.
\end{itemize}

\ejemplo

A partir de la lectura de la siguiente situación, identifica la información contenida en cada representación y propone títulos coherentes con una posible pregunta de estudio. El director de la escuela hizo una encuesta y solicita a los alumnos su colaboración para que le propongan títulos adecuados para la tabla y el gráfico y que además escriban un informe corto con el análisis de los resultados.

\fig{g03_p28_2}{Medio de transporte usado para llegar al colegio: caminando 10, bicicleta 14, bus 25, taxi 5, carro particular 12; total 66.}

% PDF p. 29
\dba{11}{Plantea y resuelve preguntas sobre la posibilidad de ocurrencia de situaciones aleatorias cotidianas y cuantifica la posibilidad de ocurrencia de eventos simples en una escala cualitativa (mayor, menor e igual).}

\evidencias

\begin{itemize}
  \item Formula y resuelve preguntas que involucran expresiones que jerarquizan la posibilidad de ocurrencia de un evento, por ejemplo: imposible, menos posible, igualmente posible, más posible, seguro.
  \item Representa los posibles resultados de una situación aleatoria simple por enumeración o usando diagramas.
  \item Asigna la posibilidad de ocurrencia de un evento de acuerdo con la escala definida.
  \item Predice la posibilidad de ocurrencia de un evento al utilizar los resultados de una situación aleatoria.
\end{itemize}

\ejemplo

Para una salida pedagógica que se realizará el día 28 de ese mes, se quiere saber si es necesario llevar impermeable. Se realiza un registro de si llueve o no durante varios días y con base en esa información se toma la decisión. Esta es la tabla que se elaboró:

\fig{g03_p29_1}{Calendario del mes con registro diario del tiempo: del 16 al 27 predominan días lluviosos (nubes), con sol los días 19 y 21; el día 28 es la salida.}

Plantea algunas ideas acerca del estado del tiempo el día de la salida, a partir de la lectura de la tabla.

